\section{Conclusions and Future Work}


We have implemented a proof-search procedure for the calculus $\Calcw$  that supports countermodel extraction.
The prover, named
\texttt{gwref}~\cite{gwrefProver}, performs a standard backward
depth-first proof search, relying on the JTabWb
engine~\cite{FerrariFF:17a}. The consistency of atomic constraints is
verified using the Choco-solver Java library~\cite{ChocoSolver:2025}.
We emphasize that the calculus $\Calcw$ and the associated
proof-search strategy can be smoothly adapted to treat the logic
$\GWCL$~\cite{FerFioRod:2025}, characterized by crisp $\GWM$-models,
by refining the notion of consistency for a set of atomic constraints
$\Gat$: a solution $\s$ to $\Gat$ must satisfy $\s(R(w,w')) \in
\{0,1\}$ for every pair of labels $w$, $w'$.  To the best of our
knowledge, the only other prover available for modal fuzzy logics is
mNiBLoS~\cite{Vidal:16}, an SMT-based solver designed for continuous
t-norm-based logics, including G\"odel--Dummett logic.  However,
neither $\GWL$ nor $\GWCL$ are supported by mNiBLoS.


The results presented in this paper are highly relevant to the study
of intermediate predicate logics.
It is well known that
modal logics can be viewed as characterizing interesting fragments of
first-order logics, in particular, the one-variable fragments of
intuitionistic predicate logic. 
From this perspective, our results have both theoretical and
practical significance.
For instance, in~\cite{BI-IGPL2016},
the authors establish Herbrand's Theorem and the
interpolation property, and introduce an alternative Skolemization
method for a wide class of intermediate intuitionistic predicate
logics enjoying the finite model property.

\tr{The results presented in this paper are highly relevant to the
  study of intermediate predicate logics. It is well known that modal
  logics can be used to characterize relevant fragments of
  intermediate first-order logics, especially certain one-variable
  fragments. From this perspective, the witnessed semantics and finite
  model property established here suggest strong connections with
  proof-theoretic properties of intermediate predicate logics, such as
  Herbrand’s Theorem, interpolation, and alternative forms of
  Skolemization studied in \cite{BI-IGPL2016} for broad classes of
  intermediate intuitionistic predicate logics with the finite model
  property. Motivated by these connections with intermediate logics,
  and in order to place $\GWL$ within the broader setting of
  intuitionistic modal logics, we plan to provide a birelational
  Kripke semantics characterization of $\GWL$, extending to this
  setting the approach developed for $\GWCL$ in
  \cite{FerFioRod:2025}. }


As future work, we plan to extend the procedure to other witnessed
G\"odel modal logics imposing additional constraints on the
accessibility relation (e.g., reflexivity, transitivity, seriality).
In contrast, treating G\"odel modal logics with non-witnessed models
seems to be considerably more challenging.




%%% Local Variables: 
%%% mode: latex
%%% TeX-master: "goedelModalLogicWitnessNonCrisp"
%%% End: 
